

\documentclass[11pt]{article}
\usepackage{amsmath,amssymb,amsthm}
\usepackage{graphicx}
\usepackage{setspace}
\usepackage[colorlinks=true,urlcolor=blue,linkcolor=blue,plainpages=false,pdfpagelabels
]{hyperref}

%\usepackage{hyperref}

\usepackage{times}

%-------------------------------------
% IOANA
%\addtolength{\textwidth}{3.5cm}
%\addtolength{\hoffset}{-2cm}
%
%\addtolength{\textheight}{3cm}
%\addtolength{\voffset}{-1cm}
%-------------------------------------

% Stanga/dreapta: 2.5cm
% Sus/jos: 3cm
\setlength{\topmargin}{-1cm}
\setlength{\textheight}{22.5cm}
\setlength{\textwidth}{16.5cm}
\setlength{\oddsidemargin}{0in}
\setlength{\evensidemargin}{0in}


\newcommand{\Z}{{\mathbb Z}}
\newcommand{\real}{{\mathbb R}}
\newcommand{\Q}{{\mathbb Q}}
\newcommand{\sra}{{\rightarrow}}
\newcommand{\Dia}{\Diamond}
\newcommand{\vsp}[1]{\vspace*{#1cm}}

\begin{document}


\begin{center}
{\Large{\bf  Curriculum vitae}}
\end{center}

\vspace{.2cm}

%\begin{minipage}[b]{0.3\linewidth}

%\noindent \includegraphics[width=0.6\textwidth]{eu.jpg}      
 
%\end{minipage}
\hspace{0.5cm}
%\begin{minipage}[b]{0.6\linewidth}
%\centering


\noindent\textbf{Ioana Leu\c{s}tean}\\
\noindent\textbf{Professor}\\
 \noindent\textbf{Department of Computer Science,\\
\noindent Faculty of Mathematics and Computer Science,\\ 
\noindent University of Bucharest, Romania.}\\
\noindent\url{https://cs.unibuc.ro/~ileustean/}\\
\noindent{email: ioana.leustean@unibuc.ro}
%\end{minipage}


\vspace{.3cm}



%\noindent
%\textbf{Personal data}
%\begin{itemize}
%    \vspace{-.3cm}
%    \item Born on 21 October 1971 in Bucharest, Romania.
%    \vspace{-.3cm}
%    \item Married, three children.
%\end{itemize}
%%----------------------------------------
%\vspace{0.2cm}
\noindent
\textbf{Education}

\noindent
[2017] Habilitation in Computer Science, University of Bucharest, Romania

Habilitation Thesis: {\em "Topics in many-valued logics"}

\noindent
[2004] Ph.D. in Mathematics, University of Bucharest, Romania

Ph.D. Thesis: {\em "Contributions to the theory of MV-algebras. MV-modules"}

Scientific Supervisor: {\em Prof. Sergiu Rudeanu}

\noindent
[1996] M.Sc. in Fundamentals of Computing, University of Bucharest, Romania

\noindent
[1995] B.A. in Mathematics, University of Bucharest, Romania

%----------------------------------------
\vspace{.2cm}
\noindent
\textbf{Positions held}

\noindent
[October 2017 - present] Professor,  University of Bucharest, Romania

\noindent
[October 2012 - September 2017] Associate Professor, University of Bucharest, Romania

\noindent
[1999 - September 2012] Assistant Professor, University of Bucharest, Romania

\noindent
[1996 - 1999] Teaching Assistant, University of Bucharest, Romania


%----------------------------------------
\vspace{.2cm}
\noindent
\textbf{Administrative positions}

\noindent
[February  2020 - February 2024] Dean, Faculty of Mathematics and Computer Science,  University of Bucharest, Romania.
%-------------------------

\vspace{.2cm}
\noindent
\textbf{Research  interests}

\noindent Current research interests are in the area of logic for specification and  verification, with a particular focus in many-sorted modal logic.   Former research belongs to the area of  {\em many-valued logic}, which I approached from various perspectives: logical, algebraic, probabilistic. 

\noindent{\em Current PhD students}: Bogan Macovei 

\noindent{\em Former PhD students}: Natalia Moanga (Ozunu)


%----------------------------------------
\vspace{.2cm}
\noindent
{\large \textbf{Awards}}

\noindent The 2008 {\em"Grigore Moisil"} Prize of the Romanian
Academy, awarded in 2010.


%----------------------------------------
\vspace{.2cm}
\noindent
{\large \textbf{Grants}} 

\noindent [October 2015 - September 2017] Grant CNCS PN-II-RU-TE-2014-4-0730
 
University of Bucharest, Romania

Project:  {\em "Modelling uncertainty in non-classical logics"} 

Grant position:  Project Director



%----------------------------------------
\vspace{.2cm}
\noindent
\textbf{Postdoctoral fellowships}



\noindent
[October 2010 - March 2013] POSDRU/89/1.5/S/58852

University of Bucharest, Romania

Project: {\em Many-Valued Logics: theoretical aspects and applications}

\noindent
[September 2007 - August 2009] "Alexander von Humboldt" Research Fellowship

Technische Universit\"{a}t Darmstadt, Germany

Project: {\em Rings, Modules and Riesz Spaces in the theory of MV-algebras}



\vspace{.2cm}
\noindent
\textbf{Past editorial activities}

\noindent  [- February, 2023] Member of the editorial board of
\href{http://www.springer.com/engineering/computational+intelligence+and+complexity/journal/500}{Soft
Computing}.


%----------------------------------------

\noindent
\textbf{Referee for peer-reviewed journals}

\noindent
Fuzzy Sets and Systems, Journal of Multiple-Valued Logic and Soft Computing, Discrete Mathematics, Order, Huston Journal of Mathematics, Journal of Pure and Applied Algebra, Algebra Universalis, Indagationes Mathematicae.

%\vspace{.2cm}
%\noindent
%\textbf{Professional memberships}
%
%\noindent
%Member of ACM (Association for Computing Machinery) and EUSFLAT (European Society for Fuzzy Logic and Technology).


\vspace{.2cm}
\noindent
\textbf{Teaching experience}

\noindent
[2009 - present] Lecturer for the following  courses:
\begin{itemize}
 \vspace{-.3cm}
    \item "Special Topics in Logic and Security" (master course, 2nd year, one module)
 \vspace{-.3cm}
    \item "Implementation of Concurrent Systems in Programming Languages" (master course, 1st year)
\vspace{-.3cm}
    \item "Declarative Programming" (undergraduate course, 3rd year)
 
    \vspace{-.3cm}
    \item "Web Techniques" (undergraduate course, 2nd year)
      
  \vspace{-.3cm}
    \item "Logic Programming" (undergraduate course, 2nd year)
    \vspace{-.3cm}
    \item "Mathematical and Computational Logic" (undergraduate course, 1st year)
\end{itemize}
%----------------------------------------


%==========================================
%\newpage
\section*{Conferences}
\vspace{0.3cm}
\noindent
{\large \textbf{Invited talks}}
\medskip

\noindent [T10] SYNASC 2020: 22nd International Symposium on Symbolic and Numeric Algorithms for Scientific
Computing (online event), 

Timisoara, Romania.

\url{https://synasc.ro/2020/}

\noindent [T9] DACS: Days of Computer Science 2016, 

Bucharest, Romania, July 3-4, 2016.

\url{http://fmi.unibuc.ro/dacs2016/}

\noindent [T8]  Soft Computing Days 2016, 

Salerno, Italy, May 23-25, 2016.

\url{http://logica.dmi.unisa.it/workshop-soft-computing-days/}

\noindent [T7] Beyond True and False: Logic, Algebra and Topology,

Florence, Italy, December 3-5, 2014 


\url{http://local.disia.unifi.it/Beyond2014/}


Title: {\em Applications of the semisimple tensor product of MV-algebras}


\noindent [T6] 35th Linz Seminar on Fuzzy Set Theory,

Linz, Austria, February 18-22, 2014.

\url{http://www.flll.jku.at/div/research/linz2014/index.html}


Title: {\em Probabilities  in \L ukasiewicz logic}


\noindent [T5] "Proof", An International Conference within the Frame of Humboldt Kollegs,

Bern, Switzerland, September 9-13, 2013.

\url{http://www.humboldt-kolleg.iam.unibe.ch/hb_main.html}


Title: {\em Linearity issues in \L ukasiewicz logic}


\noindent[T4] Anniversary Conference: Faculty of Science 150 years,

Bucharest, Romania, August 29-September 1, 2013 

\url{http://fmi.unibuc.ro/FMI-150/}

Title:{\em A propositional logic related to the Pierce-Birkhoff conjecture}


\noindent [T3 ]  "ManyVal’12",
Salerno, Italy, July 4-7, 2012.

\url{http://logica.dmi.unisa.it/manyval12/}


Title: {\em MV-modules}

\vspace{.1cm}\noindent [T2]  "Logic,
Algebra and Truth Degrees" (LATD 2010), Prague, September 7-11,
2010.

 \url{http://www.mathfuzzlog.org/latd2010/index.php}


Title: {\em Linearity issues in MV-algebras}

\vspace{.1cm}\noindent [T1] "Probability,
Uncertainty and Rationality", Pontignano, Siena, Italy, November
1-3, 2009.

 \url{http://homepage.sns.it/hosni/lori/events/pura09/}


Title: {\em Algebra and probability in Lukasiewicz logic}


%----------------------------------------

\vspace{0.3cm}
\noindent
{\large \textbf{Organization}}
\medskip

\noindent {\textbf {International workshops as main organizer}} 
\begin{itemize}
\vspace{-.3cm}
\item {\em "ManyVal 2019"}, November 1-3, 2019, Bucharest, Romania
            

\vspace{-.3cm}            
\item {\em "Algebra and Probability in Many-Valued Logics"}, May 7-9, 2009, Darmstadt, Germany 
            (financed by Alexander von Humboldt Foundation) 
\end{itemize}
\vspace{.2cm}            
\noindent {\textbf{ Special sessions at international conferences}}
\begin{itemize}
\vspace{-.3cm}
\item  {\em "Automata, logic and infinite games"} at Computability in Europe  CiE 2015,             	   	Bucharest, Romania,  June 29 – July 3, 2015 (with Dietmar Berwanger)

\vspace{-.3cm}
\item {\em "Formal Methods to Deal with Uncertainty of Many-Valued Events"} at 14th  International Conference on Information  Processing and Management of Uncertainty in Knowledge-based Systems  IPMU 2012, Catania, Italy, July 9-13, 2012 (with Tommaso Flaminio and Enrico Marchioni)
\end{itemize}

\section*{Publications}


\vspace{0.3cm}

\noindent
{\large \textbf{Papers in peer-reviewed journals}}

\noindent 
[J31] N.  Moangă, I. Leu\c stean, T.F. Şerbănuţă, Many-Sorted Hybrid Modal Languages, JLAMP 120,  Journal of Logical and Algebraic Methods in Programming 120, 100644, 2021.


\noindent
[J30] N.  Moangă, I. Leu\c stean, T.F. Şerbănuţă, A many-sorted polyadic modal logic, Fundamanenta Informaticae, 173 (2-3), 191-215,2020.

\noindent 
[J29] S. Lapenta, I. Leu\c stean, On the semisimple tensor product of MV-algebras, Fuzzy Sets and Systems 197, 140-151, 2020.

\noindent 
[J28] A. Di Nola, S. Lapenta, I. Leu\c stean, Infinitary logic and basically disconnected compact Hausdorff spaces, Journal of Logic and Computation 28(6), 1275-1292, 2018.

\noindent 
[J27] A. Di Nola, S. Lapenta, I. Leu\c stean,  An analysis of the logic of Riesz Spaces with strong unit, Annals of Pure and Applied Logic 169(3), 216-234, 2018

\noindent 
[J26] S. Lapenta, I. Leu\c stean, Notes on dvisible MV-algebras, Soft Computing, 21(21), 6213-6223, 2017.

\noindent
[J25] S. Lapenta,  I. Leu\c stean. Stochastic independence for probability MV-algebras.
      Fuzzy Sets and Systems 298, 94–206, 2016.

\noindent
[J24] S. Lapenta,  I. Leu\c stean. Scalar extensions for algebraic structures of Lukasiewicz logic. 
      Journal of Pure and Applied Algebra, 220: 1538-1553, 2016, doi:10.1016/j.jpaa.2015.09.017. 

\noindent
[J23] S. Lapenta,  I. Leu\c stean. Towards understanding the Pierce-Birkhoff conjecture via MV-algebras. 
       Fuzzy Sets and Systems,  276: 114-130, 2015.

\noindent
[J22] D. Diaconescu, I. Leu\c stean.  The Riesz MV-algebra hull of an MV-algebra.  Mathematica Slovaca,  65(4): 801–816, 2015.  Special issue in honor of Antonio Di Nola.

\noindent
[J21] D. Diaconescu, I. Leu\c stean. Mutually exclusive nuances of truth in Moisil logic.  Scientific Annals of
      Computer Science "Alexandru Ioan Cuza" Univ. of Iasi, 25:69-88, 2015.
      Special issue dedicated to Professor Sergiu Rudeanu on his 80th birthday.


\noindent 
[J20] A. Di Nola, I. Leu\c stean. \L ukasiewicz logic and Riesz spaces. Soft Computing, 18(12): 2349-2363, 2014.

\noindent
[J19] D. Diaconescu, T.  Flaminio,  I. Leu\c stean.  Lexicographic MV-algebras and lexicographic states. 
     Fuzzy Sets and Systems,  244:63-85, 2014. 

\noindent
[J18] I. Leu\c stean. Hahn-Banach theorems for MV-algebras.
 Soft Computing,   16(11): 1845-1850, 2012

\noindent
[J17] I. Leu\c stean. Metric completions of MV-algebras with states. An approach to stochastic independence. Journal of Logic and Computation, 21(3):493-508, 2011.

\noindent
[J16] I. Leu\c stean. Tensor products of probability MV-algebras.
Journal of Multiple-Valued Logic and Soft Computing, 16(3-5):405-419, 2010.

\noindent
[J15] I. Leu\c stean. A determination principle for  algebras of $n$-valued \L ukasiewicz logic. Journal of Algebra, 320: 3694-3719, 2008.

\noindent
[J14] I. Leu\c stean. $\alpha$-convergence and complete distributivity in MV-algebras. Journal of Multiple-Valued Logic and Soft Computing, 12(3-4):309-315, 2006.

\noindent
[J13] G. Georgescu, I. Leu\c stean, A. Popescu. Order convergence and distance on \L ukasiewicz-Moisil algebras. Journal of Multiple-Valued Logic and Soft Computing, 12(1-2):33-69, 2006.

\noindent
[J12] B. Gerla, I. Leu\c stean. Similarity MV-algebras. Fundamenta Informaticae, 69(3):287-300, 2006.

\noindent
[J11] I. Leu\c stean. Non-commutative \L ukasiewicz propositional logic. Archive for Mathematical Logic, 45(2):191-213, 2006.

\noindent
[J10] P. Flondor, I. Leu\c{s}tean. MV-algebras with operators: the commutative and the non-commutative case. Discrete Mathematics, 274(1-3):41-76, 2004.

\noindent
[J9] A. Di Nola, P. Flondor, I. Leu\c{s}tean. MV-modules. Journal of Algebra, 267(1):21-40, 2003.

\noindent
[J8] P. Flondor, I. Leu\c{s}tean. Tensor Products of MV-algebras. Soft Computing, 7:446-457, 2003.

\noindent
[J7] R. Ball, G. Georgescu, I. Leu\c stean. Cauchy completions of MV-algebras. Algebra Universalis, 47: 367-407, 2002.

\noindent
[J6] I. Leu\c stean. Local pseudo-MV algebras. Soft Computing, 5(5):386-395, 2001.

\noindent
[J5] G. Georgescu, I. Leu\c stean. Towards a probability theory based on Moisil logic. Soft Computing, 4(1):19-26, 2000.

\noindent
[J4] G. Georgescu, I. Leu\c stean. A representation theorem for monadic Pavelka algebras. Journal of Universal Computer Science, 6(1):105-111, 2000.

\noindent
[J3] G. Georgescu, I. Leu\c stean. Probabilities on \ Lukasiewicz-Moisil algebras. International Journal of Approximate Reasoning, 18(3-4):201-215, 1998.

\noindent
[J2] G. Georgescu, A. Iorgulescu, I. Leu\c stean. Monadic and closure MV-algebras. International Journal of Multiple-Valued Logic, 3(3):235-257, 1998.

\noindent
[J1] G. Georgescu, I. Leu\c stean. Convergence in perfect MV-algebras. Journal of Mathematical Analysis and Application, 228(1):96-111, 1998.


\vspace{.3cm}
\noindent
{\large \textbf{Journal issues as editor}}

\noindent
[JI1]  Algebra and Probability in Many-Valued Reasoning.
    Special issue. Studia Logica. 94(2), 2010. Edited by Ioana Leu\c stean and Vincenzo Marra. 

\newpage

%========================================
\vspace{0.3cm}
\noindent
{\large \textbf{Papers in  proceedings}}

\noindent
[P13] I. Leu\c stean, B. Macovei, A Hybrid Many-Sorted Modal Logic with Nominal Terms (with Ioana Leuștean), RAMiCS 2026 (accepted).

\noindent
[P12] I. Leu\c stean, B. Macovei, Łukasiewicz Logic with Actions for Neural Networks training, FROM 2025,  EPTCS 427, 2025, 44-58. 

\noindent 
[P11] A. Nicolae, P. Irofti, I. Leu\c stean, OpenBSD formal driver verification with SeL4, the 16th
Proceedings of the 16th International Conference on Security for Information Technology
and Communications , LNCS 14534, Springer, 2023

\noindent
[P10] B. Macovei,  I. Leu\c stean, DELP: Dynamic Epistemic Logic for Security Protocols, FROM 2021,  23rd International Symposium on Symbolic and Numeric Algorithms for Scientific Computing (SYNASC), 275-282, 2021.

\noindent
[P9] N.  Moangă, I. Leu\c stean, T.F. Şerbănuţă,  From Hybrid Modal Logic to Matching Logic and back, FROM 2019, EPTCS 303, 2019, 16-31. arXiv:1907.05029

\noindent
[P8] N.  Moangă, I. Leu\c stean, T.F. Şerbănuţă,  Operational semantics and program verification using many-sorted hybrid modal logic, TABLEAUX 2019, 446-476. 

\noindent
[P7] D. Diaconescu, I. Leu\c stean, Towards game semantics for nuanced logics, FUZZ-IEEE 2017, doi: 10.1109/FUZZ-IEEE.2017.8015600. 

\noindent 
[P6] A. Di Nola, B. Gerla, I. Leu\c stean. Adding Real Coefficients to Lukasiewicz Logic: An Application to Neural Networks. In:
 F. Masulli, G. Pasi, and R. Yager (Eds.), WILF 2013, LNAI 8256, pp. 77–85, 2013.



\noindent
[P5] I. Leu\c stean. State-complete Riesz MV-algebras and L-measure spaces. Proceedings of IPMU 2012, part II, CCIS, 298:226-234, 2012.

\noindent
[P4] D. Diaconescu, I. Leu\c stean, L. Petre, K. Sere, G. Stefanescu. Refinement-Preserving Translation from Event-B to Register-Voice Interactive Systems. Proceedings of IFM 2012, 221-236, 2012.

\noindent
[P3]  A. Di Nola, I. Leustean.  Riesz MV-algebras and their logic. Proceedings of EUSFLAT-LFA 2011, 140-145, 2011.

\noindent
[P2] I. Leu\c stean. The tensor PMV-algebra of an MV-algbebra. Proceedings of the 41st IEEE International Symposium for Multiple-Valued Logic, 274-276, 2011.

\noindent
[P1] B. Gerla, I. Leu\c stean. Many-valued logics and similarities. Proceedings of IPMU 2004, 477-484, 2004.

%========================================
\vspace{0.3cm}
\noindent
{\large \textbf{Contributions in handbooks and collections}}

\noindent
[C2]
S. Lapenta,  I. Leu\c stean. A general view on normal form theorems for Lukasiewicz logic with product.
      Dieter Probst and Peter Schuster (eds.), "Concepts of Proof in Mathematics, Philosophy and
      Computer Science". Ontos Mathematical Logic 6. Walter de Gruyter, Berlin, 2016. 


\noindent
[H1] A. Di Nola, I. Leustean. \L ukasiewicz logic and MV-algebras,
P. Cintula, P. Hajek, C. Noguera (eds.), Handbook of Mathematical Fuzzy Logic - volume 2,
    Studies in Logic 37, College Publications, London, 1-102,  2011.

\noindent
[C1] A. Di Nola, G. Georgescu, I. Leu\c stean. States on Perfect MV-algebras, in V. Nov\' ak, I. Perfilieva (eds.), Discovering the World With Fuzzy Logic, Stud. Fuzziness Soft Comput. 57, Physica, Heidelberg, 105-125, 2000.



\end{document}
